\begin{enumerate}
\item Đếm $n(\Omega)$ bằng quy tắc nhân, hoán vị, chỉnh hợp hoặc tổ hợp. Ví dụ: tung đồng xu $k$ lần thì $n(\Omega)=2^k$; gieo xúc xắc $k$ lần thì $n(\Omega)=6^k$; chọn $k$ phần tử từ $n$ phần tử (không thứ tự) thì $n(\Omega)=C_n^k$.
\item Đếm $n(A)$ trực tiếp bằng cách liệt kê hoặc dùng quy tắc cộng, quy tắc nhân.
\item Nếu $A$ có nhiều trường hợp hoặc có chữ "ít nhất", hãy đếm $n\left(\overline{A}\right)$ rồi tính $n(A)=n(\Omega)-n\left(\overline{A}\right)$.
\end{enumerate}
Lưu ý: dùng $A_n^k$ khi có thứ tự (xếp chỗ, phân công chức vụ, lập số) và $C_n^k$ khi không có thứ tự (chọn nhóm). Lấy có hoàn lại thì dùng quy tắc nhân $n^k$. Khi đếm theo từng trường hợp, kiểm tra để không đếm trùng.
